% v8, May 27 2026, Claude Opus 4.7

\documentclass[11pt]{article}

\usepackage[margin=1in]{geometry}
\usepackage{amsmath, amssymb, amsthm}
\usepackage{mathrsfs}
\usepackage{mathtools}
\usepackage{xcolor}
\usepackage{hyperref}
\hypersetup{colorlinks=true,
            linkcolor=blue!50!black,
            citecolor=blue!50!black,
            urlcolor=blue!50!black}

%---- math macros ----
\newcommand{\R}{\mathbb{R}}
\newcommand{\Sm}{\mathscr S}
\newcommand{\Ldual}{L^{\vee}}
\newcommand{\dd}{\mathrm d}
\newcommand{\Lag}{\Lambda}
\DeclareMathOperator{\Hess}{Hess}

%---- theorem environments ----
\newtheorem{theorem}{Theorem}
\newtheorem{proposition}{Proposition}
\newtheorem{lemma}{Lemma}
\newtheorem{corollary}{Corollary}
\theoremstyle{definition}
\newtheorem{definition}{Definition}

\title{Smoothing Polyhedral Lagrangian Surfaces}
\author{}
\date{May 27, 2026}

\begin{document}
\maketitle

\section*{Statement and Answer}

Equip $\R^4$ with coordinates $(q_1,q_2,p_1,p_2)$ and the standard symplectic
form $\omega=\dd p_1\wedge \dd q_1+\dd p_2\wedge \dd q_2$. A \emph{polyhedral
Lagrangian surface} $K\subset\R^4$ is a finite polyhedral complex, all of whose
faces are Lagrangian, that is a topological submanifold of $\R^4$.

\medskip
\noindent\textbf{Theorem.} \emph{If $K$ is a polyhedral Lagrangian surface in
which exactly $4$ faces meet at every vertex, then $K$ admits a Lagrangian
smoothing: a Hamiltonian isotopy $K_t$ of smooth Lagrangian submanifolds
parametrized by $t\in(0,1]$, extending to a topological isotopy on $[0,1]$ with
$K_0=K$.}

\medskip
\noindent The answer is \textbf{YES}.

\medskip
The proof has a local part (a normal form near each vertex and edge, from which
a smooth structure on $K$ is read off) and a global part (a fibration of
Lagrangian planes transverse to $K$, with respect to which all sufficiently
$C^1$-small ``smoothing functions'' yield genuine graphical Lagrangians). The
$4$-faces hypothesis enters in the local normal form (Lemma~\ref{lem:linear}),
where four Lagrangian quadrants force a standard symplectic basis; it is not a
Maslov/parity condition.

\section*{1. Local model at a vertex}

\begin{lemma}\label{lem:linear}
Let $\R^2\to\R^4$ be an embedding that is linear on each of the four quadrants,
with Lagrangian image $\Sigma$, and assume $\Sigma$ is not contained in a single
plane. Then there is a linear symplectic transformation of $\R^4$ carrying
$\Sigma$ to the product of the union of the positive coordinate axes in
$\R^2_{p_1q_1}$ with that in $\R^2_{p_2q_2}$.
\end{lemma}

\begin{proof}
Let $(v_1,v_2,u_1,u_2)$ be the tangent vectors at the origin to the four edges of
$\Sigma$, ordered so that cyclically adjacent vectors span the faces. The
pairings $\omega(v_i,u_i)$ cannot vanish: if some $\omega(v_i,u_i)=0$, then
$\omega$ would vanish identically on a $3$-dimensional subspace of $\R^4$, which
is impossible for a symplectic form. After swapping $v_i\leftrightarrow u_i$ if
needed we may take both pairings positive, and after rescaling we may take
$\omega(v_1,u_1)=\omega(v_2,u_2)=1$ with the cross pairings vanishing (the faces
are Lagrangian). Then $(v_1,v_2,u_1,u_2)$ is a standard symplectic basis, and the
linear map $\partial_{p_i}\mapsto v_i$, $\partial_{q_i}\mapsto u_i$ is the
required symplectomorphism, carrying $\Sigma$ to the stated product of
positive-axis unions.
\end{proof}

In the plane $\R^2_{pq}$ the symplectic pairing projects the union of the two
positive axes homeomorphically onto the dual of the antidiagonal line $p=q$.
Taking the product and applying Lemma~\ref{lem:linear}:

\begin{corollary}\label{cor:Ldual}
There is a linear Lagrangian plane $L\subset\R^4$ such that the symplectic
pairing $\R^4\to \Ldual$ restricts to a homeomorphism $\Sigma\to \Ldual$.
\end{corollary}

This equips $\Sigma$ with a smooth structure (pulled back from $\Ldual$), which
we fix henceforth. Given $L$, call a Lagrangian $\Lag\subset\R^4$
\emph{graphical} if the pairing restricts to a diffeomorphism $\Lag\cong\Ldual$.

\section*{2. Smoothing functions}

Choose a Lagrangian splitting $\R^4\cong L\oplus\Ldual$. Each of the four
quadrant-faces of $\Sigma$ is graphical over $\Ldual$, hence the graph of the
differential of a quadratic form; collecting these gives quadratic forms
$\{q_{ij}\}_{i,j\in\pm}$ on $\Ldual$, agreeing to first order along the edges.
Via $\Sigma\cong\Ldual$ (Corollary~\ref{cor:Ldual}) write $q_\Sigma$ for the
resulting $C^1$ function on $\Sigma$ that is $q_{ij}$ on each face.

\begin{definition}\label{def:smoothing}
The space $\Sm(\Sigma)$ of \emph{smoothing functions} is the set of $C^1$
functions $f:\Sigma\to\R$ such that $f+q_\Sigma$ is $C^\infty$.
\end{definition}

$\Sm(\Sigma)$ is invariant under adding smooth functions, and depends only on
$L$, not on the splitting (a different splitting changes $q_{ij}$ by a global
quadratic form, i.e.\ a smooth function). To a smoothing function $f$ we
associate the Lagrangian $\Lag_{\dd f}\subset\R^4$ that is, piecewise, the graph
of $\dd f$ over each face — equivalently the graph of $\dd(f+q_\Sigma)$ in the
splitting.

\begin{lemma}\label{lem:bijection}
$f\mapsto \Lag_{\dd f}$ is a bijection between graphical Lagrangians and
smoothing functions on $\Sigma$ modulo additive constants.
\end{lemma}

\begin{proof}
In the splitting, $\Lag_{\dd f}$ is the graph of $\dd(f+q_\Sigma)$ as a function
on $\Ldual$; graphical Lagrangians over $\Ldual$ are exactly graphs of
differentials of smooth functions, and $f+q_\Sigma$ is smooth precisely when
$f\in\Sm(\Sigma)$. The construction does not use the splitting, so the bijection
depends only on $L$.
\end{proof}

\noindent The same statements hold verbatim for an edge: a pair of Lagrangian
half-planes meeting along a line, with the analogue of
Corollary~\ref{cor:Ldual} provided by the next lemma.

\begin{lemma}\label{lem:edge}
If $\Sigma$ is a pair of linear Lagrangian half-planes meeting along a line
$\ell$, the space of Lagrangian planes $L$ for which the pairing
$\Sigma\to\Ldual$ is a homeomorphism is contractible.
\end{lemma}

\begin{proof}
Up to affine linear symplectomorphism, $\Sigma$ is the product of the real axis
in one $\R^2$-factor with the union of the positive axes in another. The pairing
$\Sigma\to\Ldual$ is a homeomorphism iff $L$ is transverse to both half-planes,
i.e.\ the symplectic reduction of $L$ along $\ell$ is a line in $\ell^\perp/\ell\cong\R^2$
meeting the interior of the positive quadrant — a contractible condition. The
space of Lagrangian lifts of such a line is itself contractible (two lifts differ
by the graph of a map $\ell'\to\ell$), giving the claim.
\end{proof}

\section*{3. Existence of small smoothing functions}

The crux is that $\Sm(\Sigma)$ is \emph{not} invariant under rescaling, so the
existence of small smoothing functions is not automatic.

\begin{lemma}\label{lem:small}
$\Sm(\Sigma)$ contains functions of arbitrarily small $C^1$-norm.
\end{lemma}

\begin{proof}
Fix a partition of unity $\sum_\sigma \chi_\sigma=1$ on $\Sigma$ indexed by the
strata, with $\chi_\sigma$ supported near $\sigma$ and $\equiv1$ on $\sigma$ away
from its boundary, of bounded $C^k$ norms. Let $\chi_\sigma^\varepsilon$ be the
composition of $\chi_\sigma$ with dilation by $1/\varepsilon$; these are
uniformly bounded with $\|\chi_\sigma^\varepsilon\|_{C^1}=O(1/\varepsilon)$.

Choose a Lagrangian plane $\Lambda_\sigma$ containing each stratum $\sigma$ and
transverse to $L$, with associated smoothing function $f_\sigma$. The first-order
agreement of the faces along shared edges means $f_\sigma$ and $f_{\sigma'}$ agree
to first order along $\sigma\cap\sigma'$. Set $f^\varepsilon:=\sum_\sigma
\chi_\sigma^\varepsilon f_\sigma$. The partition-of-unity structure makes
$f^\varepsilon$ a smoothing function, and although $\|\nabla\chi_\sigma^\varepsilon\|$
grows like $1/\varepsilon$, it is supported where $f_\sigma-f_{\sigma'}=O(\varepsilon^2)$;
by the product rule the $C^1$-norm of $f^\varepsilon$ is therefore $O(\varepsilon)$.
\end{proof}

\section*{4. Global construction: a conormal fibration}

To globalize, we replace the single plane $L$ by a smoothly varying family.

\begin{definition}\label{def:fibration}
A \emph{conormal fibration dual to $K$} is a smoothly varying family $L_z$ of
affine Lagrangian planes, $z\in K$, such that: (i) near each vertex the $L_z$ are
translates of a single plane as in Corollary~\ref{cor:Ldual}; (ii) $L_z$ varies
smoothly along edges and along faces; and (iii) in a collar of each edge the
planes in the normal direction are translates of those along the edge, with
\emph{matched} inward normal fields on the two adjacent faces (opposite vectors
under the identification $T_z\sigma\cong\Ldual_z\cong T_z\sigma'$).
\end{definition}

\begin{lemma}\label{lem:fibexists}
$K$ admits a conormal fibration which near each vertex agrees with the local
choice of Corollary~\ref{cor:Ldual}.
\end{lemma}

\begin{proof}
Lemma~\ref{lem:edge} extends the vertex choices over the edges (the space of
admissible $L_z$ along an edge is contractible). Choosing a normal vector field
to one face along each edge determines matched normals on the other; extending to
the interior of the faces is unobstructed because the space of Lagrangian planes
transverse to a given one is contractible.
\end{proof}

The fibration determines an open neighbourhood $\nu K$ of $K$ in which the fibres
$L_z$ are disjoint, hence a projection $\nu K\to K$. A Lagrangian contained in
$\nu K$ is \emph{graphical} (with respect to $L_z$) if its projection to $K$ is a
diffeomorphism. The local correspondence of Lemma~\ref{lem:bijection} globalizes:

\begin{lemma}\label{lem:graph-global}
Every graphical Lagrangian with respect to $\{L_z\}$ is the graph of a smoothing
function on $K$; and every smoothing function of sufficiently small differential
defines a graphical Lagrangian.
\end{lemma}

\begin{proof}
The correspondence is local on $K$. At a point $z$, a plane $\Ldual_z$ transverse
to $L_z$ stays transverse to nearby fibres, so a neighbourhood of $z$ in $\nu K$
is modelled on the conormal bundle of $\Ldual_z$ (Weinstein's tubular
neighbourhood theorem), and Lagrangians there are graphs of closed $1$-forms,
i.e.\ differentials of smoothing functions.
\end{proof}

\begin{lemma}\label{lem:small-global}
There exist smoothing functions for $K$ of arbitrarily small $C^1$-norm.
\end{lemma}

\begin{proof}
Take a partition of unity $\sum_\alpha\rho_\alpha=1$ on $K$ indexed by strata,
with $\rho_\alpha$ supported in the open star of $\alpha$. Lemma~\ref{lem:small}
gives smoothing functions $f_\alpha$ of arbitrarily small $C^1$-norm on each open
star; then $\sum_\alpha\rho_\alpha f_\alpha$ is a global smoothing function of
small $C^1$-norm.
\end{proof}

\section*{5. Assembly of the smoothing}

\begin{proof}[Proof of the Theorem]
By Lemma~\ref{lem:fibexists} fix a conormal fibration $\{L_z\}$ with disjoint
fibres on a neighbourhood $\nu K$. By Lemma~\ref{lem:small-global} choose a
sequence of smoothing functions $f_i\to 0$ in $C^1$ with $\dd f_i$ small enough
that, by Lemma~\ref{lem:graph-global}, each $\Lag_{\dd f_i}\subset\nu K$ is a
smooth embedded graphical Lagrangian. These $\Lag_{\dd f_i}$ are all isotopic to
$K$ by a piecewise-smooth isotopy and converge to $K$ in the Hausdorff topology
as $i\to\infty$ (equivalently as the $C^1$-norm $\to0$).

Because $\Lag_{\dd f_i}$ and $\Lag_{\dd f_{i+1}}$ are graphs of differentials of
smooth functions over the common base $K$ (via the fibration), the linear
interpolation of those functions produces a smooth Hamiltonian path of graphical
Lagrangians connecting them. Concatenating these paths and reparametrizing yields
a Hamiltonian isotopy $K_t$ ($t\in(0,1]$) of smooth Lagrangians with $K_t\to K$
as $t\to0$, extending to a topological isotopy on $[0,1]$ with $K_0=K$. This is
the desired Lagrangian smoothing.
\end{proof}

\section*{Role of the $4$-faces-per-vertex hypothesis}

The hypothesis is used exactly once, in Lemma~\ref{lem:linear}: four Lagrangian
quadrants meeting at a vertex, with edge tangents $(v_1,v_2,u_1,u_2)$, force
$\omega(v_i,u_i)\ne0$ and hence a standard symplectic basis, which is what makes
the symplectic pairing $\Sigma\to\Ldual$ a homeomorphism
(Corollary~\ref{cor:Ldual}) and supplies the smooth structure used throughout.
With a different number of faces the local model changes and this normal form is
unavailable; the mechanism is the linear symplectic geometry of the vertex, not a
Maslov-class parity condition.

\section*{Remarks}

\begin{enumerate}\setlength\itemsep{4pt}
\item \textbf{Why smoothing functions, not Legendrian fillings.} The smoothing is
produced directly as graphs over $K$ via the conormal fibration, so no
holomorphic-curve or contact-fillability input is needed; the only analytic
ingredients are Weinstein's neighbourhood theorem and a partition-of-unity
estimate.
\item \textbf{The rescaling subtlety.} $\Sm(\Sigma)$ is closed under adding
smooth functions but not under scalar multiplication, so small smoothing
functions must be built by hand (Lemma~\ref{lem:small}); this is the one place a
genuine estimate is required.
\item \textbf{Matched normals.} Property (iii) of Definition~\ref{def:fibration}
ensures the two smooth structures on $K$ — from Corollary~\ref{cor:Ldual} and from
the edge collars — agree, so the global graph construction is consistent across
edges.
\item \textbf{Generalization.} The argument is local-to-global and uses only
linear symplectic normal forms plus partition-of-unity gluing, so it adapts to
other ambient symplectic manifolds once the vertex normal form is established.
\end{enumerate}

\begin{thebibliography}{9}
\bibitem{Hirsch} M.~W.~Hirsch.
  \emph{Differential Topology}. Graduate Texts in Mathematics, Springer, 1976.
\bibitem{McDuffSalamon} D.~McDuff and D.~Salamon.
  \emph{Introduction to Symplectic Topology}, 3rd ed.,
  Oxford University Press, 2017.
\end{thebibliography}

\end{document}